\section{相关工作与差异}

区域或场群风电预测需要同时刻画时间变化和场站间空间相关性。早期深度学习方法利用 CNN-LSTM 建模风速时序特征 \cite{ref1}，也有工作将风机嵌入规则网格，使卷积网络能够处理空间分布信息 \cite{ref2}。围绕场群时空关联的研究进一步将风电场站间的地理距离、气象相关性、功率相关性和趋势相似性纳入建模，已在层级图建模、动态图建模、趋势感知、空间信息集成和持续学习等场群预测任务中取得进展 \cite{ref3,ref4,ref5,ref6,ref7,ref8,ref9,ref10,ref11,ref12}。这些研究说明，多场站功率状态能够提供有价值的空间耦合信息，但其重点通常是历史功率动态或场群空间结构建模，而不是在未来气象主驱动下进行状态条件化融合。

未来气象预报是多步风电预测中的主要外生信息。已有方法将 NWP 与历史功率进行多模态融合，或围绕 NWP 预报误差、时间滞后、多源时空特征和空间相关性进行建模 \cite{ref13,ref14,ref15,ref16,ref17,ref18}。对于网格化气象场，也有研究使用 3D CNN、区域深度学习框架或概率预测框架直接编码气象数据立方体 \cite{ref19,ref20,ref21}。本文采用 Pangu-Weather 作为未来气象来源 \cite{ref22}，并将其视为预测窗口内的主驱动信息。

Transformer 及其变体为多步时间序列预测提供了并行建模基础。Informer、PatchTST、Crossformer 和 Non-stationary Transformers 分别从长序列注意力、时间片段表示、跨维依赖和非平稳性角度改进序列预测 \cite{ref23,ref24,ref25,ref26}。这些方法说明，注意力机制适合处理多步预测中的长时依赖，但并不直接回答业务化风电预测中未来气象与预测起点运行状态如何分工的问题。

上述两类研究分别强调场站功率时空关系和未来气象驱动，但在业务预测协议下二者并不对称。未来气象在整个预测窗口内可用，而真实功率只在预测窗口前可用。因此，本文并不追求让两类输入以相同方式进入模型，而是让功率状态在气象读取阶段起条件化作用。这一点与常见的多模态分支融合有所区别。

与现有特征拼接、并行分支或对称注意力融合不同，本文关注 memory-level 的非对称融合。query-based decoding 和可查询记忆结构已在结构化输出和通用感知模型中得到应用 \cite{ref27,ref28,ref29}。本文借鉴这一形式，但不进行语言 prompt learning 或外部检索增强；dynamic power context tokens 完全由当前样本的场站级功率动态生成，并用于条件化 future queries 对气象记忆的读取过程。
